\subsection{The residual Bernoulli test in the higher-rank argument}

The auxiliary character used in the deformation argument is odd. Its
first Bernoulli number can be a unit even when the target value
\(L(-e,\psi)\) vanishes by parity. We give the auxiliary unit with
the precise conductor threshold, and the resulting cohomological
implication in both character directions.

\begin{theorem}[Unit for the auxiliary two-power twist]\label{famsel:unit}
Let \(\ell\) be prime, \(d\ge1\), \(Q=\ell^d\), and
\(a=v_2(Q-1)\ge2\). Let \(f\) be odd, \(\phi\) primitive of conductor
\(f\) and order dividing \(Q-1\), and \(\xi\) primitive of conductor
\(M=2^n\). Take the character values in \(\overline{\mathbf Q}_\ell\).
If \(M\ge f\,4^a\) and \(\phi(-1)\xi(-1)=-1\), then
\(B_{1,\phi\xi}\) is an \(\ell\)-adic unit. This includes \(\ell\mid f\).
\end{theorem}
\begin{proof}
Put \(m=n-a\) and \(s=f2^a\). Then \(m\ge a\ge2\) and
\(1<s\le2^m\). Extend characters by zero at nonunits of their modulus.
For \(u\in\mathbf Z/M\mathbf Z\) set
\[
 A(u)=\sum_{\substack{0\le b<fM\\ b\equiv u\bmod M}} b\phi(b),
 \qquad S=\sum_{b\bmod f}\phi(b),\qquad
 g(u)=\frac{A(u)-A(0)}{fM}.
\]
The sum of the unweighted \(\phi(b)\) in each progression is \(S\),
since \(M\) is coprime to \(f\). If \(t\ge0\), shifting the letters
of this progression by \(ft\) gives
\[
 g(u)-g(u-ft)=tS/M-
       \sum_{\substack{0\le b<ft\\b\equiv u\bmod M}}\phi(b)
       \qquad(0\le ft\le fM).
\]
Indeed the terms which wrap past \(fM\) lose \(fM\) in the weighted
sum, and \(\phi\) has period \(f\). Choose \(0\le t<M\) with
\(ft\equiv u\bmod M\). This formula proves \(g(u)\) integral at \(\ell\)
and \(\overline{g(u)}\in\mathbf F_Q\), since \(g(0)=0\), \(\ell\)
is odd and the nonzero values of \(\phi\) reduce to \(\mathbf F_Q\).
In particular no division by \(f\) is needed on reduction.
Write \([u]\) for the least nonnegative representative. At \(t=2^a\)
we obtain
\[
 \bar g(u)-\bar g(u-s)=c-{\bf1}_{[u]<s}\,\bar\phi([u]),
 \qquad c=\overline{2^a S/M}.
\]
The character \(\xi\) has sum zero. Grouping the defining first
Bernoulli sum by progressions gives
\(B_{1,\phi\xi}=\sum_u\xi(u)g(u)\), as \(\phi\xi\) is nontrivial.

Suppose the reduction of this last sum vanishes. For every unit \(y\)
we also have \(\sum_u\bar\xi(u)\bar g(uy)=0\), by the substitution
\(uy=v\). Apply normalized finite-field trace from
\(\mathbf F_Q(\mu_{2^n})\) to \(\mathbf F_Q\). Its degree is a power
of two, invertible in characteristic \(\ell\). The trace fixes roots
of order dividing \(2^a\) and kills every root of order \(2^j\)
with \(j>a\). Indeed \(Q^{2^{j-a-1}}\equiv1+2^{j-1}\pmod{2^j}\)
by successive squaring of \(Q\equiv1\bmod4\). The corresponding
Frobenius automorphism takes such a root to its negative. Thus
\[
 \sum_{u\in D}\bar\xi(u)\bar g(uy)=0,\qquad
 D=\{u\in(\mathbf Z/M\mathbf Z)^\times:\bar\xi(u)^{2^a}=1\}.
\]
Every \(u\in D\) is congruent to \(1\) or \(-1\) modulo \(2^m\).
In fact the units are \(\{\pm1\}\times\langle5\rangle\) with
\(5\) of order \(2^{n-2}\), by \(v_2(5^{2^j}-1)=j+2\).
Primitivity gives \(\bar\xi(1+M/2)=-1\), so \(\bar\xi\) is faithful
on \(\langle5\rangle\). Writing \(u=\pm5^i\), we have
\(2^{m-2}\mid i\), which gives the claimed congruence.
Let \(\epsilon(u)=1\) on the first congruence class and \(-1\)
on the second. Multiplying by \(1+M/2\) preserves the two classes
and \(D\) and negates \(\bar\xi\). Hence
\(\sum_D\epsilon(u)\bar\xi(u)=0\).

Subtract the identities for \(y=1-s\) and \(y=1\), taking the latter
minus the former. Since \(2^a\mid s\), we have \(u(1-s)=u-s\)
in the positive class and \(u(1-s)=u+s\) in the negative class.
In the first class the only boundary contribution \([u]<s\) occurs
at \(u=1\), where it contributes \(-1\). In the second class use the
shift formula at \(u+s\). The inequality \([u+s]<s\) forces
\([u+s]=s-1\), because \(s\le2^m\) and \(u+s\equiv s-1\bmod 2^m\).
Here \(u=-1\), and its contribution is
\(\bar\xi(-1)\bar\phi(s-1)=-1\). Both exceptional units belong to
\(D\). The resulting equality is
\(0=c\sum_D\epsilon(u)\bar\xi(u)-2=-2\), impossible at odd \(\ell\).
\end{proof}
\src{src/Development/AutomorphicArithmetic/AutomorphicVarieties/LevelCohomology/LxPartial.lean}{15}{progression difference},
\src{src/Development/AutomorphicArithmetic/AutomorphicVarieties/LevelCohomology/LxCore.lean}{65}{two boundary contributions}, and
\src{src/Development/AutomorphicArithmetic/AutomorphicVarieties/LevelCohomology/LemmaX.lean}{78}{exact conductor threshold}.

At the threshold itself take \(\ell=5,\ d=f=1,\ M=16\), and the
primitive character with \(\xi(5)=i,\ \xi(-1)=-1\).
Pairing the residues \(b=1,5,9,13\) with \(16-b\) gives
\[
 16B_{1,\xi}=(2-16)+(10-16)i-(18-16)-(26-16)i=-16(1+i).
\]
The norm of \(B_{1,\xi}\) is \(2\), a unit at either place over five.

\begin{proposition}[The residual classes annihilated by the unit]\label{famsel:stick}
Let \(\eta\ne1\) be an integral Dirichlet character modulo \(D_0\),
of order prime to an odd prime \(\ell\). Set \(\gamma=\bar\eta\)
as a Galois character and suppose it differs from the mod-\(\ell\)
cyclotomic character. If \(B_{1,\eta^{-1}}\) is a unit at \(\ell\),
every continuous cocycle \(z:G_{\mathbf Q}\to k(\gamma)\), for a
discrete residue field \(k\) containing its values, is a coboundary
under the following local condition. On the intersection of each
finite inertia group with \(\ker\gamma\), the cocycle vanishes on a
subgroup of finite index prime to \(\ell\).
In Theorem~\ref{famsel:unit} this applies to both
\(\eta=\phi\xi\) and \(\eta=(\phi\xi)^{-1}\).
\end{proposition}
\begin{proof}
On \(\ker\gamma\) the cocycle is an additive homomorphism, so the
local vanishing extends to the whole inertia intersection: a nonzero
finite image would have order divisible by \(\ell\), dividing the
given index. On \(G_L\), \(L=\mathbf Q(\mu_{D_0})\), it therefore
factors through \(\mathrm{Cl}(L)\) by global class field theory.
Real inertia has order at most two and is killed as well. If
\(\sigma_b(\zeta)=\zeta^b\), conjugating a group element in \(G_L\)
by a lift of \(\sigma_b\) multiplies the restricted homomorphism by
\(\bar\eta(b)\), by the cocycle identity and equivariance of reciprocity.

Choose a positive odd integer \(c\) prime to \(D_0\) with
\(c-\eta(c)\) a unit at \(\ell\). Indeed choose a Galois element on
which \(\gamma\) and the residual cyclotomic character differ and
lift its unit class modulo \(\operatorname{lcm}(2,\ell,D_0)\).
The classical Stickelberger theorem says that
\[
 (c-\sigma_c)\Theta\in\mathbf Z[\operatorname{Gal}(L/\mathbf Q)]
 \quad\text{annihilates }\mathrm{Cl}(L),\qquad
 \Theta=D_0^{-1}\!\!\sum_{b\in(\mathbf Z/D_0\mathbf Z)^\times}
                  [b]\,\sigma_b^{-1}.
\]
This is the integral annihilator for the full cyclotomic field
\cite{stickelberger1890}. Evaluating the integral group-ring element
at \(\eta\) in characteristic zero gives
\((c-\eta(c))B_{1,\eta^{-1}}\), since \(\eta^{-1}\ne1\).
Reduction of this scalar is nonzero even if \(\ell\mid D_0\).
Thus the restricted homomorphism is zero. The cocycle descends to
the abelian group \(\operatorname{Gal}(L/\mathbf Q)\). Choose \(h\)
with \(\gamma(h)\ne1\), possible since reduction is faithful on the
prime-to-\(\ell\) roots. Comparing \(z(hg)\) and \(z(gh)\) gives
\(z(g)=(\gamma(g)-1)z(h)/(\gamma(h)-1)\).

In the application the character and its inverse both have odd parity
and order prime to \(\ell\). Primitivity gives \(\xi(1+M/2)=-1\).
Choose a unit congruent to \(1+M/2\) modulo \(M\) and to \(1\) modulo
\(\ell f\), using coprimality. On a Galois lift both residual characters
equal \(-1\), whereas the cyclotomic character is \(1\).
Theorem~\ref{famsel:unit} applies to the inverse pair of characters
as well. This proves both instances.
\end{proof}
\src{src/Development/AutomorphicArithmetic/AutomorphicVarieties/LevelCohomology/ResidualBernoulliNonvanishing.lean}{174}{residual cocycle criterion},
\src{src/Development/AutomorphicArithmetic/AutomorphicVarieties/LevelCohomology/ChrInvNeg.lean}{197}{both auxiliary characters}, and
\src{src/Development/AutomorphicArithmetic/AutomorphicVarieties/ModularityArgument/ThresholdArguments/DenseThresholdCases/ConductorAtMapInjectiveCoeff.lean}{132}{use in the two higher-rank counts}.
These two statements discharge the residual Bernoulli premise of the
rank-four counting theorem. Its other premises include integral
deformation finiteness, rigidity and an automorphic count. The cocycles
in Proposition~\ref{famsel:stick} have residual coefficients in an
auxiliary odd character, whereas the required lengths count all
congruence depths in the original coefficient module.

\subsection{Integral vanishing for varying nontrivial characters}

\begin{proposition}\label{famsel:class}
Let \(\ell\) be odd, \(O\) a coefficient DVR finite over
\(\mathbf Z_\ell\), with uniformizer \(\pi\), and \(\psi\) a Dirichlet
character of order prime to \(\ell\) valued in \(O^\times\).
Put \(\alpha=\psi^{-1}\chi_\ell^{-e}\),
\(A=(\operatorname{Frac}O/O)(\alpha)\) and
\(K=\overline{\mathbf Q}^{\ker\bar\alpha}\), where the bar denotes
reduction modulo \(\pi\). Suppose \(\bar\alpha(D_\ell)\ne1\).
Put \(s=[\mathbf F_\ell(\bar\alpha(G_{\mathbf Q})):\mathbf F_\ell]\).
If the order of the group
\[
        \operatorname{Cl}(K)/\langle[\mathfrak l]:\mathfrak l\mid\ell\rangle
\]
is not divisible by \(\ell^s\), the entire strict group
\(S_\ell(\mathbf Q,A)\) is zero. Its normalized length at every host
level \(\ell N\) is then zero.
\end{proposition}
\begin{proof}
Each continuous cocycle into the discrete module \(A\) has finite
image, so each class is killed by a power of \(\pi\). Take a strict
class killed by \(\pi\), and subtract a coboundary so that its
representative \(z\) vanishes on \(D_\ell\).
We have \(\pi z(g)=g b-b\) for some \(b\in A\). Thus
\(b\in A^{D_\ell}=0\), since some element acts by a scalar whose
difference from one is a unit. The values of \(z\) lie in \(A[\pi]\),
so its restriction to \(G_K\) is a homomorphism. It vanishes on
every decomposition group above \(\ell\) in \(G_K\), by the cocycle
conjugation formula and the trivial residual action of \(G_K\).
For \(p\ne\ell\), inertia acts trivially through \(\chi_\ell\) and
reduction is injective on the prime-to-\(\ell\) image of \(\psi\).
Thus \(I_p\cap G_K\) acts trivially on the full \(A\). On that
subgroup the restriction of \(z\) is a coboundary by strictness,
hence is zero. Inertia at infinity has order at most two and has no
image in \(A[\pi]\). By global reciprocity this homomorphism
factors through the displayed quotient of the ordinary ideal class
group. Its image in \(A[\pi]\) is stable under every residual scalar,
by the cocycle conjugation formula. It is an additive subgroup of
the residue field, hence a vector space over the field generated by
these scalars. A nonzero image would have order divisible by \(\ell^s\).
Thus it vanishes. Inflation--restriction leaves a class of
\(\operatorname{Gal}(K/\mathbf Q)\) in \(A[\pi]\). This Galois
group has order dividing \(\#(O/\pi)^\times\), so averaging kills
the class. Vanishing of the \(\pi\)-torsion proves the claim at every depth.
\end{proof}

In particular, let the primitive conductor \(f\) be prime to \(\ell\),
let \(\psi\) be even and nontrivial, and take any positive multiple
\(e\) of \(\ell-1\). If \(\psi(\ell)\ne1\), the local hypothesis holds.
The field \(K\) is the totally real field cut out by \(\psi\), independent
of \(e\). Writing \(d=\operatorname{ord}(\psi)\), one explicit sufficient
condition for vanishing, uniform in the character and the twist, is
\[
        \ell^{\operatorname{ord}_d(\ell)}>
              f^{d/2}\bigl(1+\tfrac d2\log f\bigr)^{d-1}.
\]
Indeed \(|\Delta_K|\le f^d\), by the discriminant formula in the tower
\(\mathbf Q(\mu_f)/K/\mathbf Q\) and the elementary cyclotomic bound
\(|\Delta_{\mathbf Q(\mu_f)}|\le f^{\varphi(f)}\).
Minkowski puts a representative of every ideal class at norm at most
\(B=\sqrt{|\Delta_K|}\). At a rational prime \(p\) the coefficients of
the Euler factor for counting ideals are bounded by those of
\((1-p^{-s})^{-d}\), since there are at most \(d\) prime ideals and
their residue degrees are positive. Counting ordered products of \(d\)
positive integers with product at most \(B\) therefore bounds \(h_K\)
by \(B(1+\log B)^{d-1}\), using the harmonic sum \(d-1\) times.
Here the residual character still has order \(d\), so the stated
inequality gives \(h_K<\ell^s\). This extends the odd-order test of
Proposition~\ref{newcoh:oddzero} to either parity of the character
order and accounts for the required splitting above \(\ell\).

Apply the class-group criterion of Proposition~\ref{famsel:class} directly. The quadratic character modulo \(5\) at \(\ell=7\)
gives zero for every \(e=6j\), and either order-three character
modulo \(7\) at \(\ell=5\) gives zero for every \(e=4j\), \(j\ge1\).
The Frobenius residues are \(2\bmod5\) and \(5\not\equiv\pm1\bmod7\).
The first field is \(\mathbf Q(\sqrt5)\), with Minkowski bound
\(\sqrt5/2<2\). The second is \(\mathbf Q(\mu_7)^+\).
The polynomial \(X^3+X^2-2X-1\) of \(\zeta_7+\zeta_7^{-1}\)
has discriminant \(49\), so its Minkowski bound is at most
\(6\cdot7/27<2\). Both class groups are trivial.
Whenever these criteria hold for an admitted host character in case II
or III, they give the required polynomial inequality at all the indicated
twists. Characters whose residual action is locally trivial at \(\ell\),
and those for which the class-group test fails, still enter the counts.
\src{src/Development/AutomorphicArithmetic/AutomorphicVarieties/HeightAndVanishing/Counts.lean}{65}{Case II}
and \src{src/Development/AutomorphicArithmetic/AutomorphicVarieties/HeightAndVanishing/Counts.lean}{67}{case III}.

\begin{proposition}[A host-level class test allowing \(\ell\)-power inertia]
\label{famsel:hostinertia}
Let \(\ell\) be odd, \(N\) a host modulus, and \(\psi\) any character
modulo \(N\) with values in \(O^\times\). Use \(A,\bar\alpha,K,s\)
as in Proposition~\ref{famsel:class}. Let
\[
 T=\{p\mid\ell N:\bar\alpha(D_p)\ne1\}.
\]
Suppose \(\ell\in T\), and every \(p\ne\ell\) with
\(\ell\mid\#\psi(I_p)\) belongs to \(T\).
If \(\ell^s\) does not divide the order of
\(\operatorname{Cl}(K)/\langle[\mathfrak p]:\mathfrak p\mid p,\ p\in T\rangle\),
the full strict group at level \(\ell N\) vanishes.
\end{proposition}
\begin{proof}
Normalize a \(\pi\)-torsion class to a cocycle \(z\) vanishing on
\(D_\ell\) as in Proposition~\ref{famsel:class}. Then
\(z\) takes values in \(A[\pi]\). For \(p\in T\), the host condition
gives \(z|_{D_p}=\delta b_p\) in \(A\). Since \(\pi z=0\),
\(\pi b_p\in A^{D_p}=0\). Thus \(b_p\in A[\pi]\) and
\(z|_{D_p\cap G_K}=0\). At every other \(p\ne\ell\), the order of
\(\psi(I_p)\) is prime to \(\ell\), so reduction is injective on
that group, and \(I_p\cap G_K\) acts trivially on \(A\). Strictness
on inertia again makes \(z\) zero there. The homomorphism
\(z|_{G_K}\) factors through the displayed quotient by reciprocity.
The image-size and inflation--restriction arguments of
Proposition~\ref{famsel:class} apply verbatim, since
\(\operatorname{Gal}(K/\mathbf Q)\) has order prime to \(\ell\)
regardless of the order of the original \(\psi\).
\end{proof}

\src{src/Development/AutomorphicArithmetic/AutomorphicVarieties/CharacterLevelBounds/ShaFamMono.lean}{30}{strict length at level \(\ell N\)} and
\src{src/Development/AutomorphicArithmetic/AutomorphicVarieties/GaloisLibrary/SelmerTransport/ConjStableShaFam.lean}{94}{decomposition conditions at the primes of the level}.

For a fixed exceptional prime, take \(\ell=3\) and either primitive
order-six character \(\psi\) modulo \(7\). Let \(7\mid N\) and inflate
\(\psi\) to \(N\). For every positive odd \(e\) its reduction after
the inverse Tate twist is the product of the quadratic characters of
discriminants \(-7\) and \(-3\): reduction kills the cubic part,
and inversion fixes both quadratic factors. Hence
\(K=\mathbf Q(\sqrt{21})\); its residual character is nontrivial on
inertia at both \(3\) and \(7\). The only prime away from \(3\)
where \(\psi\) has inertia of order divisible by \(3\) is \(7\).
Minkowski puts each ideal class at norm at most \(\sqrt{21}/2<3\).
The integer \((1+\sqrt{21})/2\) has polynomial \(X^2-X-5\),
irreducible modulo \(2\), so \(2\) is inert and no ideal has norm
two. Thus \(h_K=1\), and Proposition~\ref{famsel:hostinertia} gives
zero at every congruence depth. With either sextic character as the
fixed target, its inflation occurs on each of its hosts at all odd
twists, by the defining kernel of the host. For odd \(e\ge3\) these
are case II: \(3\) is exceptional and the prime-to-\(3\) part away
from \(3\) has order two. This includes arbitrarily high twists with
\(3\mid\operatorname{ord}(\psi)\).
\src{src/Development/AutomorphicArithmetic/AutomorphicVarieties/HeightAndVanishing/DirichletPrimeToLPartMulLPowerPart.lean}{179}{host kernel and parity test},
\src{src/Development/AutomorphicArithmetic/AutomorphicVarieties/HeightAndVanishing/DirichletPrimeToLPartMulLPowerPart.lean}{212}{case II}, and
\src{src/Development/AutomorphicArithmetic/AutomorphicVarieties/StmtFields.lean}{66}{the fixed exceptional set}.

\begin{proposition}[Characters of prime-power conductor]\label{famsel:primary}
Let \(\ell\) be odd and suppose \(\ell\nmid h_{\mathbf Q(\zeta_\ell)^+}\).
For every Dirichlet character \(\psi\) of conductor dividing a power
of \(\ell\) and every \(e\ge1\) with \(\psi(-1)=(-1)^e\), the strict
group \(S_\ell(\mathbf Q,A)\) vanishes, also on every host to which
\(\psi\) inflates.
\end{proposition}
\begin{proof}
The full scalar \(\alpha\) factors through the \(\ell\)-power
cyclotomic tower. This tower is totally ramified at \(\ell\), so
\(\alpha(D_\ell)=\alpha(G_{\mathbf Q})\). For a class killed by
\(\pi\), normalize \(z|_{D_\ell}=0\) and write \(\pi z=\delta b\).
Then \(b\in A^{D_\ell}=A^{G_{\mathbf Q}}\), hence \(\pi z=0\);
this works even when \(\bar\alpha=1\).
The prime-to-\(\ell\) character \(\bar\alpha\) factors through
\((\mathbf Z/\ell)^\times\), and parity makes its field \(K\)
a subfield of \(L=\mathbf Q(\zeta_\ell)^+\).
Extension of ideals embeds the \(\ell\)-part of \(\mathrm{Cl}(K)\)
in \(\mathrm{Cl}(L)\), since its composite with norm is raising to
\([L:K]\), a unit at \(\ell\).
The homomorphism \(z|_{G_K}\) is zero on inertia away from \(\ell\)
by strictness and is zero on decomposition groups over \(\ell\)
by normalization and conjugation. Ordinary class field theory now
kills it, since \(\ell\nmid h_K\). Inflation--restriction and
\(\ell\nmid[K:\mathbf Q]\) kill the original class. As before,
vanishing of the \(\pi\)-torsion suffices.
\end{proof}

In fact, for the sextic target modulo \(7\), every power of the
target has zero strict group at both exceptional primes \(3\) and
\(7\), at all twists of its admitted parity. At \(7\) this follows
from Proposition~\ref{famsel:primary} and \(h_{\mathbf Q(\zeta_7)^+}=1\).
At \(3\) the odd powers, including the quadratic character, all
have the residual field \(\mathbf Q(\sqrt{21})\) at odd twists, so
Proposition~\ref{famsel:hostinertia} applies as above. For an even
power and an even twist, \(\bar\alpha=1\). Inertia at \(3\) maps
onto the full \(3\)-adic cyclotomic image, and a Frobenius at \(3\)
generates the image of the unramified character \(\psi\) of order
one or three, since \(3\not\equiv\pm1\bmod7\). Thus
\(\alpha(D_3)=\alpha(G_{\mathbf Q})\), and the normalization just
proved makes \(z\) a homomorphism of exponent dividing three, zero
on \(D_3\). It is unramified away from \(7\) by strictness and
is tame at \(7\) by its exponent. By reciprocity it factors through
the ordinary ray group modulo \(7\) of \(\mathbf Q\), which is
\((\mathbf Z/7)^\times/\{\pm1\}\). The class of \(3\) generates
this group and maps to zero, proving the claim. This treats the
inflated target powers on each host. Other host characters can have
additional conductor factors.

\begin{proposition}[A unit test at a completely split prime]\label{famsel:splitray}
Let \(\psi\) be a nontrivial even character of order prime to an odd
prime \(\ell\). Suppose \(\ell\) splits completely in its character
field \(K=\overline{\mathbf Q}^{\ker\psi}\), and \(\ell\nmid h_K\).
At the embeddings \(w:K\hookrightarrow\mathbf Q_\ell\) put
\[
 q_w(u)=\frac{w(u)^{\ell-1}-1}{\ell}\pmod\ell
 \qquad(u\in\mathcal O_K^\times).
\]
If the vectors \((q_w(u))_w\) span the coordinate-sum-zero hyperplane
over \(\mathbf F_\ell\), then \(S_\ell(\mathbf Q,A)=0\) for every
positive multiple \(e\) of \(\ell-1\), also after imposing any host
conditions. Here \(\bar\alpha\) is trivial on \(D_\ell\).
\end{proposition}
\begin{proof}
A strict class killed by \(\pi\) lifts to a cocycle \(z\) in
\(A[\pi]\) by the exact sequence for multiplication by \(\pi\);
no local condition on this lift at \(\ell\) is needed. The residual
action is \(\bar\psi^{-1}\). On \(G_K\) the lift is a homomorphism
unramified away from \(\ell\), since every such inertia group there
acts trivially on \(A\). Its image has exponent \(\ell\). At each
\(w\mid\ell\), local reciprocity therefore kills
\(1+\ell^2\mathbf Z_\ell=(1+\ell\mathbf Z_\ell)^\ell\), by the
\(\ell\)-adic logarithm. Real inertia has order two. Thus the
homomorphism factors through the ray class group of modulus
\(\ell^2\mathcal O_K\) modulo \(\ell\)-th powers.

The ray class exact sequence and \(\ell\nmid h_K\) identify this
quotient with a quotient of
\((\mathcal O_K/\ell^2)^\times/
((\mathcal O_K/\ell^2)^\times)^\ell\simeq\bigoplus_w\mathbf F_\ell\)
by the unit vectors above. The displayed isomorphism is the
coordinatewise Fermat quotient. Its kernel is the subgroup of
\(\ell\)-th powers, and it is surjective on each factor, as is seen
on \(1+\ell a\). The vectors sum to zero since the norm of a unit is
\(\pm1\). Under the span hypothesis the ray class quotient has trivial
Galois action: the coordinate-sum map is invariant under permutation.
The restriction of \(z\) is equivariant with target of nontrivial
scalar character \(\bar\psi^{-1}\), hence must vanish. Finally
\([K:\mathbf Q]=\operatorname{ord}(\psi)\) is invertible on \(A[\pi]\),
so inflation--restriction kills the class there and in \(A\).
Every cohomology class with coefficients in \(A\) is killed by a power of \(\pi\), which proves the assertion.
For the ray class sequence see
\href{https://math.mit.edu/classes/18.785/2021fa/LectureNotes21.pdf\#page=5}{Sutherland, Theorem 21.8}.
\end{proof}

\paragraph{A full case III progression, including the split local action.}
For either cubic character modulo \(7\) and \(\ell=29\), the strict
group vanishes for \emph{every} positive \(e\) with \(28\mid2e\).
First take \(28\mid e\). The field \(K=\mathbf Q(\zeta_7)^+\) has
class number one by the Minkowski computation above, and \(29\equiv1\bmod7\)
splits completely. The unit \(u=\zeta_7+\zeta_7^{-1}\) has polynomial
\(X^3+X^2-2X-1\); \(u^2-2\) is a conjugate. Modulo \(29^2=841\)
the three roots, ordered cyclically under \(r\mapsto r^2-2\), are
\(409,761,511\). Direct exponentiation gives
\[
 (q_w(u))_w=(20,16,22),\qquad (q_w(u^2-2))_w=(16,22,20).
\]
Both vectors have coordinate sum \(0\bmod29\), and their first
two columns have determinant \(20\cdot22-16^2\equiv10\bmod29\).
Proposition~\ref{famsel:splitray} applies at every multiple of \(28\).

For \(e\equiv14\bmod28\), Proposition~\ref{newcoh:oddzero}(2)
reduces the question to \(29^2\nmid h_{K(\sqrt{29})}\), since
\(\operatorname{ord}_3(29)=2\). The compositum is totally real of
degree six and discriminant at most \(49^2\,29^3\), by the order
obtained by tensoring integral bases. Its Minkowski bound is
at most \((6!/6^6)49\,29^{3/2}<120\).
The primes \(2,3,5\) all have degree three in \(K\), by their
classes modulo \(7\) up to sign. Also \(2\) is inert in
\(\mathbf Q(\sqrt{29})\), so its degree in the compositum is six.
Consequently no prime ideal has norm below seven and the only ones
of norm at most ten lie over \(7\). There are at most two of these,
since \(7\) is totally ramified in \(K\).
Every ideal of norm \(<120\) is the unit ideal, a prime ideal, or a
product of two prime ideals. There are at most \(6\cdot60\) prime
ideals of norm \(<120\), since at most 60 rational primes are below
120. In a product the smaller norm is at most ten and the other is
at most \(120/7<18\), giving at most \(2\cdot6\cdot17\) products.
Hence \(h_{K(\sqrt{29})}\le1+360+204<29^2\), as required.

Taking a cubic character modulo \(7\) as the target gives
\(p_K=11,\ p_1=13\) in the exceptional set defined below. The class
number of \(\mathbf Q(\sqrt{-11})\) is one: its Minkowski bound
\((2/\pi)\sqrt{11}<3\) and \(2\) is inert. Thus \(29\notin B_\chi\).
The characters inflate to every host for this target. The two
computations exhaust the twists in case III at this prime, uniformly
in their size and in the host.

\begin{corollary}[Case III on two families of targets]\label{famsel:hostdensity}
For every target of two-power order, including the zeta target,
case III is empty on every one of its hosts. For either primitive
cubic target modulo \(7\), there is an explicit set of coefficient
primes of relative natural density \(1/2\) among primes on which
every case III strict group vanishes for all sufficiently large
primes in the set, simultaneously in the host and the twist.
\end{corollary}
\begin{proof}
In the host group of Proposition~\ref{host:art}, the kernel of
the projection to the cyclic target image is a two-group. Every
odd-order character is therefore inflated from that cyclic image.
If the target order is a power of two, only the trivial character
is possible, whereas case III requires a nontrivial one. This
applies in particular to \(\psi=1\). There are still case II tests
on the zeta hosts: for \(\ell=3\) and each odd \(e\ge3\) in the
window take \(J_d\), where \(d>3\) is the auxiliary prime.
It is a host character, is odd since \(d\equiv3\bmod4\), and its
prime-to-three part away from three has order two.

For a cubic target the only nontrivial odd-order host characters
are the two cubic characters modulo \(7\). Their common field is
\(K=\mathbf Q(\zeta_7)^+\), of class number one. Take primes
\[
    \ell\bmod 7\in\{2,3,4,5\},\qquad
    \ell\bmod 12\in\{5,7,11\}.
\]
Both characters are nontrivial at \(\ell\). If \(\ell-1\mid e\),
Proposition~\ref{newcoh:oddzero}(1) gives vanishing. For
\(\ell\equiv3\bmod4\), admissible \(e\) is even, so
\(\ell-1\mid 2e\) already gives \(\ell-1\mid e\).
In the only remaining residue class, \(\ell\equiv5\bmod12\),
the order of \(\ell\) modulo three is \(2>3/2\).
Proposition~\ref{newcoh:oddzero}(2) gives vanishing for the other
twists for all sufficiently large such \(\ell\), with a threshold
depending only on \(K\). Imposing the further host local conditions
preserves zero. There are \(4\cdot3=12\) such reduced classes
out of \(\varphi(84)=24\) by the Chinese remainder theorem.
The prime number theorem for this fixed modulus gives their
relative natural density. Omitting the finitely many exceptional
or small primes preserves that density.
\end{proof}

\subsection{An effective bound simultaneous in the conductor, twist and prime}

\begin{theorem}\label{famsel:rate}
For \(N,e\ge1\), sum the normalized strict lengths of
Proposition~\ref{newcoh:fullsupport} over every prime and any set of
distinct characters of modulus \(N\) at each prime. The coefficient
fields may vary. With \(P=16Ne\) one has
\[
 \sum_{\ell,\psi}\frac{\log\ell}{e_{\ell,\psi}}
       \operatorname{length}_{O_{\ell,\psi}}S_{\ell,\psi}
       \ \le\ \exp(P^{102}).
\]
\end{theorem}
\begin{proof}
Put \(M=4N,\ F=\mathbf Q(\mu_M),\ d=[F:\mathbf Q],\ j=2e\)
and \(H=\max(d+1,2j+2)\). The group \(K_j(\mathcal O_F)\) is finite.
For its primary parts with \(\ell>H\), Bayer-Fluckiger, Emery
and Houriet's Theorem~1.1~\cite{bayerEmeryHouriet} bounds the sum of
their log orders by
\[
 B=(2j+1)^{71j^4d^3}\, d^{293j^5d^5}\,
             |\operatorname{disc}F|^{528j^5d^4}.
\]
Indeed the subgroup generated by elements of order at most \(H\)
has zero projection to these primary parts. We only use the
large-prime consequence of their theorem.

For \(\ell\le H\), let \(F'=F(\mu_\ell)\).
Its degree over \(F\) divides \(\ell-1\). The extension of the two
integer schemes after inverting \(\ell\) is finite \'{e}tale:
adjoining \(\mu_\ell\) is unramified outside \(\ell\).
On \'{e}tale \(H^2\) with the Tate twist \(e+1\), trace composed
with pullback is multiplication by \([F':F]\). It is invertible
at \(\ell\), by the finite-level trace identity and passage to the
inverse limit. Via the even Quillen--Lichtenbaum comparison in
Theorem~\ref{newcoh:lemmaK}, the primary group over \(F\) thus embeds
in the one over \(F'\). The proof of that theorem bounds the order of
the entire primary \(K\)-group, by \eqref{newcoh:transfer} and its
subsequent estimates. The hypotheses on roots of unity hold over
\(F'\) also at two, since \(4\mid M\). We have \(d'=[F':\mathbf Q]\le dH\)
and \(|\operatorname{disc}F'|\le(MH)^{d'}\).
Indeed \(F'\) has conductor dividing \(\operatorname{lcm}(M,\ell)\);
the derivative of \(X^c-1\) at a primitive \(c\)-th root is divisible
in the integers by \(\Phi_c'\) there. Taking norms gives the usual
discriminant bound \(c^{\varphi(c)}\) for the cyclotomic field,
using \(\mathbf Z[\zeta_c]=\mathcal O_{\mathbf Q(\mu_c)}\).
It follows that
\[
 \log\#K_j(\mathcal O_F)
 \le B+H\exp\!\left((jdH)^{32}(1+dH)\log(2MH)\right).
\]
Here \(d,j,H,M\le P\), \(2j+1\le P\), and
\(\log|\operatorname{disc}F|\le d\log M\). Consequently
\(\log B\le 892 P^{10}\log P\). The exponent in the second term
is at most \(6P^{98}\log P\), using \(P\ge16\).
Each term on the right is at most \(\exp(P^{100})\), including
the factor \(H\). Thus \(\log\#K_j(\mathcal O_F)\le\exp(P^{101})\).
Proposition~\ref{newcoh:fullsupport} now bounds the required sum by
\(\varphi(N)(\exp(P^{101})+d\log(4Nde))\le\exp(P^{102})\).
\end{proof}

\begin{corollary}[Part of the window budget]\label{famsel:subwindow}
Fix \(\varepsilon>0\). At a common host level \(N\), retain, for each
positive integer twist \(e\), any distinct characters of primitive
conductor \(f\) satisfying \(16fe\le x^{\varepsilon/204}\).
At all sufficiently large \(x\), their total normalized strict length
at level \(\ell N\), summed over the primes, twists and characters,
is at most \(\exp(x^\varepsilon/2)\). The threshold depends only on
\(\varepsilon\), independently of the host and coefficient fields.
\end{corollary}
\begin{proof}
At fixed \(N\) an inducer modulo \(f\mid N\) has a unique inflation.
A strict class at level \(\ell N\) satisfies in particular all the
conditions at level \(\ell f\). Group by \(f,e\) and apply
Theorem~\ref{famsel:rate} at \(f\). Each pair contributes at most
\(\exp(x^{\varepsilon/2})\). Bounding each coordinate separately gives
at most \(x^{\varepsilon/102}\) pairs. The sum is therefore at most
\(x^{\varepsilon/102}\exp(x^{\varepsilon/2})\), whose logarithm
\(x^{\varepsilon/2}+(\varepsilon/102)\log x\) is at most
\(x^\varepsilon/2\) for sufficiently large \(x\).
\end{proof}

\subsection{The complementary twists and the field budget}\label{famsel:complement}
Use \(v\) for the page scale, \(t=\lfloor v^{401}\rfloor\), and \(N\)
for its host conductor. The scale of the final comparison is
\(x=\kappa v\) with \(\kappa\) fixed. Already \(f=1,\ e=\lfloor v\rfloor\)
eventually violates \(16fe\le v^{\varepsilon/204}\) in the active range
\(0<\varepsilon<1/52\). We record which additional estimates enter
at that same host.

Fix the primitive target \(\chi\), of order \(o\) and conductor \(m\).
Let \(p_K\) be the least prime \(p\ge7\), \(p\equiv3\bmod4\), prime
to \(m\), let \(h_* = h_{\mathbf Q(\sqrt{-p_K})}\), and let \(p_1\)
be the least prime congruent to \(1\) modulo the odd part of \(o\)
and prime to \(2mp_K\). The fixed exceptional set is
\[
 B_\chi=\{3\}\cup\{\ell\text{ odd prime}:\ell\mid om p_K h_*p_1\}.
\]
At odd \(\ell\) use the host characters \(\psi\) with
\(\psi(-1)=(-1)^e\), \(1\le e<t\). Case II requires \(e\ge2\), even
order of the prime-to-\(\ell\) part of \(\psi\) with its local
\(\ell\)-component removed, and
\(\ell\in B_\chi\) or \(\ell-1\mid2e\). Case III requires
\(\ell\notin B_\chi,\ \ell-1\mid2e\), and nontrivial \(\psi\) of
odd order. On every host whose two auxiliary primes exceed \(2t+1\)
and \(m,o,p_K,h_*,p_1\), the imported counts give a finite normalized
length \(w_{\ell,e,\psi}=(\log\ell)\operatorname{length}_O
S_{\ell N}(\mathbf Q,A)/e(O/\mathbf Z_\ell)\) at most
\((\ell N e)^{C}\) in both cases. One exponent \(C\), the larger of
the two source exponents, works for every such host and choice of
coefficient rings for this target.
\src{src/Development/AutomorphicArithmetic/AutomorphicVarieties/StmtFields.lean}{59}{auxiliaries and exceptional primes},
\src{src/Development/AutomorphicArithmetic/AutomorphicVarieties/HeightAndVanishing/DirichletPrimeToLPartMulLPowerPart.lean}{185}{host size and six cases}, and
\src{src/Development/AutomorphicArithmetic/AutomorphicVarieties/HeightAndVanishing/Counts.lean}{60}{uniform character bounds}.

Let \(\mathcal S\) consist of all character entries with
\(16fe\le v^{\varepsilon/204}\), and \(\mathcal R\) of the cases II
and III outside \(\mathcal S\), deleting any entries proved zero
above. For \(L=\max(2t,\max B_\chi)\), every entry of \(\mathcal R\)
has \(\ell\le L\), since \(\ell-1\mid2e\) implies \(\ell\le2t-1\).
There are at most \(N\) characters per prime and twist. Applying the
explicit estimate in the proof of Corollary~\ref{famsel:subwindow}
only to the first sum therefore gives
\[
 \sum_{\mathcal S}w_{\ell,e,\psi}+\sum_{\mathcal R}w_{\ell,e,\psi}
 \le v^{\varepsilon/102}\exp(v^{\varepsilon/2})
       + LNt\,(LNt)^{C}.
\]
On the admitted page \(\log(Nt)\le 2v^{\varepsilon/4}\).
Since \(t\le v^{401}\), the logarithm of each summand is
\(o(v^\varepsilon)\), uniformly on these hosts. The sum fits in
\(\exp(v^\varepsilon/2)\) for sufficiently large \(v\).
Only the second summand uses the imported bounds in cases II and III.
\src{src/Development/AutomorphicArithmetic/AutomorphicVarieties/HeightAndVanishing/SumOut.lean}{167}{related finite-prime count outside the exceptional set} and
\src{src/Development/AutomorphicArithmetic/AutomorphicVarieties/CMFieldExponentialBounds/LogConductorMulHeightLeOfPageN.lean}{6}{uniform conductor bound}.

For clarity, the full sum over the CM field \(E\) has further terms.
Away from \(B_\chi\) and two the remaining cases are I
(\(\ell-1\nmid2e\)) and Z (the trivial character with
\(\ell-1\mid2e\)). For I the imported estimate sums over any finite
set of its primes \emph{before} bounding by \((Nt)^{C_1}\) for each
\(e\), once both auxiliary primes exceed two additional test primes
depending only on \(\chi\). Proposition~\ref{newcoh:zero} gives Z.
Inside \(B_\chi\), case IV handles \(e\ge2\) with the indicated
prime-to-\(\ell\) part of odd order via the fixed cyclotomic layer,
and the two estimates for case V handle \(e=1\).
Using the field-to-character comparison and the fixed-layer bounds,
the source obtains the odd-prime plus contribution
\((Nt)^{C'}+C'B\), where \(B=\exp(v^{4\varepsilon/5})\).
The same window bounds the \(2\)-part by \(B\log2\); the odd-prime
minus estimate is at most \(K\exp(8v^{\varepsilon/4})\).
Consequently for every finite prime set \(S\), on that field,
\[
 \begin{aligned}
 &\sum_{\ell\in S}\log\ell\sum_{1\le e<t}
      \operatorname{length}_{\mathbf Z_\ell}
      S_\ell(E,(\mathbf Q_\ell/\mathbf Z_\ell)(-e))\\
 &\qquad\le B\log2+K\exp(8v^{\varepsilon/4})+(Nt)^{C'}+C'B
       \le \exp(v^\varepsilon/2).
 \end{aligned}
\]
Here \(C',K\) are independent of \(v\) and of the admitted host for
fixed \(\chi\). The last inequality follows by taking logarithms
and using the bound on \(\log(Nt)\). The comparisons on \(E\), the
case I sum and the remaining nontrivial character bounds are used at
the indicated source theorems. The preceding character inequality
identifies the part of their input supplied by the independent estimates.
\src{src/Development/AutomorphicArithmetic/AutomorphicVarieties/HeightAndVanishing/Counts.lean}{53}{case I with both test primes},
\src{src/Development/AutomorphicArithmetic/AutomorphicVarieties/HeightAndVanishing/LayerLe.lean}{16}{fixed layer and field comparisons},
\src{src/Development/AutomorphicArithmetic/AutomorphicVarieties/HeightAndVanishing/JointWinOfCounts.lean}{16}{assembly of the counts}, and
\src{src/Development/AutomorphicArithmetic/AutomorphicVarieties/CMFieldExponentialBounds/AdmissibleOfPageN.lean}{110}{dyadic, minus and plus contributions on one field}.
